\documentclass[11pt,aspectratio=169]{ctexbeamer}

\usetheme{default}
\setbeamertemplate{navigation symbols}{}

\usepackage{amsmath}
\usepackage{booktabs}

\definecolor{DeepBlue}{RGB}{25,74,122}
\definecolor{SoftInk}{RGB}{53,59,66}
\definecolor{PaperBg}{RGB}{249,248,245}
\definecolor{PanelBlue}{RGB}{236,243,249}

\setbeamercolor{normal text}{fg=SoftInk,bg=PaperBg}
\setbeamercolor{structure}{fg=DeepBlue}
\setbeamercolor{frametitle}{fg=DeepBlue,bg=PaperBg}
\setbeamercolor{title}{fg=DeepBlue}
\setbeamercolor{block title}{fg=DeepBlue,bg=PanelBlue}
\setbeamercolor{block body}{fg=SoftInk,bg=white}

\newcommand{\Core}[1]{
  \begin{block}{核心结论}
    #1
  \end{block}
}

\newcommand{\LectureMeta}[2]{
  \title{大模型推理系统与实践}
  \subtitle{#1}
  \author{张书豪}
  \institute{华中科技大学计算机学院}
  \date{2027 · #2}
}

\newcommand{\PaperRef}[3]{
  \item \textbf{#1}（#2）\hfill\href{#3}{原文链接}
}

\newcommand{\ResearchSources}[1]{
  \begin{frame}{本讲参考资料}
    \small
    \begin{itemize}
      #1
    \end{itemize}
    \vspace{0.3em}
    \textcolor{SoftInk}{课件中的论文数字均指作者在特定模型、硬件和工作负载下报告的结果，不代表无条件性能。}
  \end{frame}
}

\newcommand{\InternalCaseNote}{
  \small\textcolor{SoftInk}{说明：本页为课程研究案例或教学假设，不代表已发表论文结论；只有在实验包给出模型、硬件、工作负载和原始记录时，数字才可作为结果引用。}
}

% 14 讲细分版的统一参考资料入口。论文数字的解释边界见课程内容基线。

\newcommand{\LectureOneSources}{\ResearchSources{
  \PaperRef{vLLM: Easy, Fast, and Cheap LLM Serving with PagedAttention}{SOSP 2023}{https://doi.org/10.1145/3600006.3613165}
  \PaperRef{SGLang: Efficient Execution of Structured Language Model Programs}{arXiv:2312.07104}{https://arxiv.org/abs/2312.07104}
  \PaperRef{nano-vllm-hust 教学代码}{课程仓库}{https://github.com/intellistream/intro-to-llm-inference-engines}
}}

\newcommand{\LectureTwoSources}{\ResearchSources{
  \PaperRef{vLLM / PagedAttention}{SOSP 2023}{https://doi.org/10.1145/3600006.3613165}
  \PaperRef{Sarathi-Serve}{OSDI 2024}{https://www.usenix.org/conference/osdi24/presentation/agrawal}
  \PaperRef{DistServe}{OSDI 2024}{https://www.usenix.org/conference/osdi24/presentation/zhong-yinmin}
}}

\newcommand{\LectureThreeSources}{\ResearchSources{
  \PaperRef{vLLM / PagedAttention}{SOSP 2023}{https://doi.org/10.1145/3600006.3613165}
  \PaperRef{DistServe}{OSDI 2024}{https://www.usenix.org/conference/osdi24/presentation/zhong-yinmin}
  \PaperRef{Llumnix}{OSDI 2024}{https://www.usenix.org/conference/osdi24/presentation/sun-biao}
}}

\newcommand{\LectureFourSources}{\ResearchSources{
  \PaperRef{Orca}{OSDI 2022}{https://www.usenix.org/conference/osdi22/presentation/yu}
  \PaperRef{Sarathi-Serve}{OSDI 2024}{https://www.usenix.org/conference/osdi24/presentation/agrawal}
  \PaperRef{FastServe}{NSDI 2026}{https://www.usenix.org/conference/nsdi26/presentation/wu-bingyang}
}}

\newcommand{\LectureFiveSources}{\ResearchSources{
  \PaperRef{vLLM / PagedAttention}{SOSP 2023}{https://doi.org/10.1145/3600006.3613165}
  \PaperRef{vAttention}{arXiv:2405.04437}{https://arxiv.org/abs/2405.04437}
  \PaperRef{Rethinking KV Cache Compression}{MLSys 2025}{https://proceedings.mlsys.org/paper_files/paper/2025/hash/26289c647c6828e862e271ca3c490486-Abstract-Conference.html}
  \PaperRef{Mooncake}{FAST 2025}{https://www.usenix.org/conference/fast25/presentation/qin}
}}

\newcommand{\LectureSixSources}{\ResearchSources{
  \PaperRef{DistServe}{OSDI 2024}{https://www.usenix.org/conference/osdi24/presentation/zhong-yinmin}
  \PaperRef{RAGCache}{arXiv:2404.12457}{https://arxiv.org/abs/2404.12457}
  \PaperRef{Mooncake}{FAST 2025}{https://www.usenix.org/conference/fast25/presentation/qin}
  \PaperRef{NEO}{MLSys 2025}{https://proceedings.mlsys.org/paper_files/paper/2025/hash/66a026c0d17040889b50f0dfa650e5e0-Abstract-Conference.html}
}}

\newcommand{\LectureSevenSources}{\ResearchSources{
  \PaperRef{vLLM / PagedAttention}{SOSP 2023}{https://doi.org/10.1145/3600006.3613165}
  \PaperRef{SGLang}{arXiv:2312.07104}{https://arxiv.org/abs/2312.07104}
  \PaperRef{DistServe}{OSDI 2024}{https://www.usenix.org/conference/osdi24/presentation/zhong-yinmin}
  \PaperRef{Mooncake}{FAST 2025}{https://www.usenix.org/conference/fast25/presentation/qin}
}}

\newcommand{\LectureEightSources}{\ResearchSources{
  \PaperRef{FlashAttention}{NeurIPS 2022}{https://arxiv.org/abs/2205.14135}
  \PaperRef{FlexGen}{ICML 2023}{https://proceedings.mlr.press/v202/sheng23a.html}
  \PaperRef{PowerInfer}{SOSP 2024 / arXiv:2312.12456}{https://arxiv.org/abs/2312.12456}
  \PaperRef{FlashInfer}{MLSys 2025}{https://proceedings.mlsys.org/paper_files/paper/2025/hash/dbf02b21d77409a2db30e56866a8ab3a-Abstract-Conference.html}
}}

\newcommand{\LectureNineSources}{\ResearchSources{
  \PaperRef{DEEPSERVE}{USENIX ATC 2025}{https://www.usenix.org/conference/atc25/presentation/hu-junhao}
  \PaperRef{QServe}{MLSys 2025}{https://proceedings.mlsys.org/paper_files/paper/2025/hash/fbe2b2f74a2ece8070d8fb073717bda6-Abstract-Conference.html}
  \PaperRef{vLLM-HUST}{开源维护仓库}{https://github.com/vLLM-HUST/vllm-hust}
}}

\newcommand{\LectureTenSources}{\ResearchSources{
  \PaperRef{vLLM / PagedAttention}{SOSP 2023}{https://doi.org/10.1145/3600006.3613165}
  \PaperRef{Revisiting Disaggregated LLM Serving}{2025 preprint, arXiv:2601.08833}{https://arxiv.org/abs/2601.08833}
  \PaperRef{Rethinking KV Cache Compression}{MLSys 2025}{https://proceedings.mlsys.org/paper_files/paper/2025/hash/26289c647c6828e862e271ca3c490486-Abstract-Conference.html}
}}

\newcommand{\LectureElevenSources}{\ResearchSources{
  \PaperRef{MLPerf Inference}{MLCommons benchmark}{https://mlcommons.org/benchmarks/inference-datacenter/}
  \PaperRef{ACM Artifact Review and Badging}{ACM policy}{https://www.acm.org/publications/policies/artifact-review-and-badging-current}
  \PaperRef{Revisiting Disaggregated LLM Serving}{2025 preprint, arXiv:2601.08833}{https://arxiv.org/abs/2601.08833}
}}

\newcommand{\LectureTwelveSources}{\ResearchSources{
  \PaperRef{nano-vllm-hust}{课程教学代码}{https://github.com/intellistream/intro-to-llm-inference-engines}
  \PaperRef{vLLM-HUST}{HUST 维护的 vLLM fork}{https://github.com/vLLM-HUST/vllm-hust}
  \PaperRef{上游 vLLM}{开源项目}{https://github.com/vllm-project/vllm}
}}

\newcommand{\LectureThirteenSources}{\ResearchSources{
  \PaperRef{课程内容基线与实验材料}{课程仓库}{https://github.com/intellistream/intro-to-llm-inference-engines}
  \PaperRef{DEEPSERVE}{USENIX ATC 2025}{https://www.usenix.org/conference/atc25/presentation/hu-junhao}
  \PaperRef{GhostServe}{MLSys 2026}{https://proceedings.mlsys.org/paper_files/paper/2026/hash/d99e8e80a6c41e148db686918dd7eab3-Abstract-Conference.html}
}}

\newcommand{\LectureFourteenSources}{\ResearchSources{
  \PaperRef{课程完整材料与后续更新}{课程仓库}{https://github.com/intellistream/intro-to-llm-inference-engines}
  \PaperRef{SpecGen}{MLSys 2026}{https://proceedings.mlsys.org/paper_files/paper/2026/hash/66a026c0d17040889b50f0dfa650e5e0-Abstract-Conference.html}
  \PaperRef{GhostServe}{MLSys 2026}{https://proceedings.mlsys.org/paper_files/paper/2026/hash/d99e8e80a6c41e148db686918dd7eab3-Abstract-Conference.html}
}}



\LectureMeta{第 7 讲 推理系统架构}{研究生课程课件}

\begin{document}

\begin{frame}[plain]
  \titlepage
\end{frame}

\begin{frame}{教学目标}
  \Core{本讲从架构层面理解现代推理 runtime，并建立论文机制与真实代码边界之间的对应关系。}
  \begin{itemize}
    \item 区分控制面与数据面。
    \item 识别调度器、执行器、KV 管理器和通信层。
    \item 理解为什么日志、脚本与回归也是系统组成部分。
  \end{itemize}
\end{frame}

\begin{frame}{典型推理系统的部件}
  \Core{一个现代 serving runtime 至少包括请求入口、调度器、执行器、缓存层和监控层。}
  \begin{itemize}
    \item API server 与 tokenizer
    \item scheduler 与 admission control
    \item model runner 与 kernel/runtime 层
    \item KV manager、通信层与 metrics/logging
  \end{itemize}
\end{frame}

\begin{frame}{为什么架构图不能只画模块名}
  \Core{好的架构理解不仅要知道“有哪些模块”，还要知道请求、状态和控制信号如何在模块间流动。}
  \begin{itemize}
    \item 请求路径回答“谁先被调用”。
    \item 状态路径回答“哪些对象长期存在”。
    \item 控制路径回答“谁在决定系统行为”。
  \end{itemize}
\end{frame}

\begin{frame}{控制面与数据面}
  \Core{控制面决定做什么，数据面决定怎么做。}
  \begin{itemize}
    \item 控制面关注准入、路由、优先级、配额和反馈。
    \item 数据面关注 token 计算、状态驻留、kernel 执行和通信。
    \item 不区分两者，问题定义就容易混乱。
  \end{itemize}
\end{frame}

\begin{frame}{控制面典型问题}
  \Core{控制面的核心是资源治理和服务目标兑现。}
  \begin{itemize}
    \item 什么时候接收新请求。
    \item 怎样分配配额和优先级。
    \item 什么时候扩容、限流或改变路由。
  \end{itemize}
\end{frame}

\begin{frame}{数据面典型问题}
  \Core{数据面的核心是执行路径、状态驻留与通信组织。}
  \begin{itemize}
    \item 哪些 token 被算、以什么批次被算。
    \item KV 和其他状态如何分配、迁移和释放。
    \item kernel、图执行和并行通信如何被组织。
  \end{itemize}
\end{frame}

\begin{frame}{为什么要进入真实代码库}
  \Core{只有进入代码，才能真正理解论文机制在系统中落在哪些接口和模块上。}
  \begin{itemize}
    \item 请求生命周期在哪里开始与结束。
    \item KV block 在哪里分配和回收。
    \item 指标输出如何与论文图表对应。
  \end{itemize}
\end{frame}

\begin{frame}{源码穿插：用 nano-vllm-hust 建立最小架构图}
  \Core{\texttt{\detokenize{nano-vllm-hust}} 可以把请求入口、调度、KV 管理和执行路径落到一套足够小的代码结构上。}
  \begin{itemize}
    \item 请求入口：\texttt{\detokenize{nanovllm/engine/llm_engine.py}}
    \item 调度与状态推进：\texttt{\detokenize{nanovllm/engine/scheduler.py}}
    \item KV 块管理：\texttt{\detokenize{nanovllm/engine/block_manager.py}}
    \item 执行与 KV 准备：\texttt{\detokenize{nanovllm/engine/model_runner.py}}
  \end{itemize}
\end{frame}

\begin{frame}{源码穿插：从架构图到代码导读任务}
  \Core{以 \texttt{\detokenize{nano-vllm-hust}} 为例，代码导读可以先收束到一条最小主路径。}
  \begin{itemize}
    \item 从 \texttt{\detokenize{LLMEngine.generate()}} 进入全局驱动循环。
    \item 从 \texttt{\detokenize{Scheduler.schedule()}} 观察控制面如何选择当前执行对象。
    \item 从 \texttt{\detokenize{ModelRunner.run()}} 观察数据面如何准备输入、block table 与采样输出。
  \end{itemize}
\end{frame}

\begin{frame}{代码导读的最低要求}
  \Core{一次合格的代码导读，应能够将关键函数与系统职责清晰对应起来。}
  \begin{itemize}
    \item 哪个函数负责接收请求。
    \item 哪个循环负责组织批次。
    \item 哪个模块负责状态分配和统计输出。
  \end{itemize}
\end{frame}

\begin{frame}{辅助模块的重要性}
  \Core{脚本、日志与回归流程共同决定系统结果能否被复现和比较。}
  \begin{itemize}
    \item benchmark 配置决定实验边界。
    \item 日志与指标决定你是否真的理解系统行为。
    \item 回归脚本决定结果是否可持续维护。
  \end{itemize}
\end{frame}

\begin{frame}{课堂练习}
  \Core{课堂练习可先做一次最小代码导读，建立架构与实现之间的映射关系。}
  \begin{itemize}
    \item 定位请求入口。
    \item 定位调度循环。
    \item 定位 KV 分配和回收逻辑。
    \item 定位 metrics 输出点。
  \end{itemize}
\end{frame}

\begin{frame}{讲解：系统架构课为什么必须把三条路径画出来}
  \Core{好的架构分析不只是列模块名，而是同时说明请求路径、状态路径和观测路径。}
  \begin{itemize}
    \item 请求路径回答“系统怎么跑起来”。
    \item 状态路径回答“哪些对象在不同模块之间被持有、更新、回收”。
    \item 观测路径回答“我们靠什么知道系统真的按预期工作”。
  \end{itemize}
\end{frame}

\begin{frame}{讲解：为什么日志、benchmark 和回归也属于架构}
  \Core{如果只把 runtime 主循环画进架构图，课程最后很容易出现“会讲代码，不会验证系统”的断裂。}
  \begin{itemize}
    \item benchmark 决定 workload 如何进入系统，结果如何被汇总。
    \item 日志和指标决定问题是否可诊断，机制是否可解释。
    \item 回归脚本决定改动能否长期维护，因此它们不是外围工具，而是系统可持续性的组成部分。
  \end{itemize}
\end{frame}

\begin{frame}{例子：为什么控制面和数据面一定要分开讲（1）}
  \Core{系统架构课的关键不是记模块名，而是知道谁负责决策、谁负责执行、谁负责携带状态。}
\begin{block}{控制面}
    接收请求、维护队列、决定调度策略、记录指标和做资源治理。
  \end{block}
\end{frame}

\begin{frame}{例子：为什么控制面和数据面一定要分开讲（2）}
  \Core{系统架构课的关键不是记模块名，而是知道谁负责决策、谁负责执行、谁负责携带状态。}
\begin{block}{数据面}
    真正执行 tokenizer、prefill、decode、KV 访问和结果回传。
  \end{block}
\end{frame}

\begin{frame}{例子：为什么控制面和数据面一定要分开讲（3）}
  \Core{系统架构课的关键不是记模块名，而是知道谁负责决策、谁负责执行、谁负责携带状态。}
\begin{block}{课程结论}
    如果控制面和数据面在讲授中混在一起，学生很难看出系统边界和可改控制点。
  \end{block}
\end{frame}

\begin{frame}{例子：分布式拓扑为什么会改变同一个问题的答案（1）}
  \Core{单机上看起来只是“多开几个请求”的问题，进入分布式拓扑后会变成通信和协调问题。}
\begin{block}{单卡}
    主要矛盾是调度、KV 驻留和本地显存组织。
  \end{block}
\end{frame}

\begin{frame}{例子：分布式拓扑为什么会改变同一个问题的答案（2）}
  \Core{单机上看起来只是“多开几个请求”的问题，进入分布式拓扑后会变成通信和协调问题。}
\begin{block}{多卡 TP / PP / DP}
    同一个 batch 决策会影响跨卡同步、负载倾斜和阶段间流水稳定性。
  \end{block}
\end{frame}

\begin{frame}{例子：分布式拓扑为什么会改变同一个问题的答案（3）}
  \Core{单机上看起来只是“多开几个请求”的问题，进入分布式拓扑后会变成通信和协调问题。}
\begin{block}{教学意义}
    架构课要让学生意识到：系统模块关系会随着部署拓扑变化而重排。
  \end{block}
\end{frame}

\begin{frame}{代码例子：从 LLMEngine 到 ModelRunner 的调用链（1）}
  \Core{系统架构课最适合补的代码页，是跨组件调用链而不是单函数细节。}
\begin{block}{调用链草图}
    {\ttfamily\small
    \detokenize{LLM.generate -> LLMEngine.generate}\par
    \detokenize{LLMEngine.step -> Scheduler.schedule}\par
    \detokenize{LLMEngine.step -> ModelRunner.execute_model}\par
    }
  \end{block}
\end{frame}

\begin{frame}{代码例子：从 LLMEngine 到 ModelRunner 的调用链（2）}
  \Core{系统架构课最适合补的代码页，是跨组件调用链而不是单函数细节。}
\begin{block}{课堂应追问}
    这一链路上，哪个位置属于控制面，哪个位置已经进入数据面？
  \end{block}
\end{frame}

\begin{frame}{代码例子：从 LLMEngine 到 ModelRunner 的调用链（3）}
  \Core{系统架构课最适合补的代码页，是跨组件调用链而不是单函数细节。}
\begin{block}{系统含义}
    架构分析的第一步不是画框图，而是把真实调用关系翻译成模块职责图。
  \end{block}
\end{frame}

\begin{frame}{代码例子：为什么可观测性也属于架构的一部分（1）}
  \Core{没有日志、指标和 trace，系统架构图只是静态想象，无法支持真实调试和研究。}
\begin{block}{最小可观测性接口}
    {\ttfamily\small
    \detokenize{logger.info(stage, request_id, batch_size)}\par
    \detokenize{metrics.observe("ttft_ms", value)}\par
    \detokenize{trace.annotate("decode_step_begin")}\par
    }
  \end{block}
\end{frame}

\begin{frame}{代码例子：为什么可观测性也属于架构的一部分（2）}
  \Core{没有日志、指标和 trace，系统架构图只是静态想象，无法支持真实调试和研究。}
\begin{block}{课堂应追问}
    哪些观测点属于控制面信号，哪些属于执行面信号？
  \end{block}
\end{frame}

\begin{frame}{代码例子：为什么可观测性也属于架构的一部分（3）}
  \Core{没有日志、指标和 trace，系统架构图只是静态想象，无法支持真实调试和研究。}
\begin{block}{结论}
    真正可维护的架构，必须把观测路径和主执行路径一起设计。
  \end{block}
\end{frame}

\begin{frame}{例子：一个坏架构图通常坏在哪里（1）}
  \Core{学生最常见的问题是画出很多框，但无法说明数据、状态和决策谁在流动。}
\begin{block}{坏图一}
    只列模块名，不说明调用方向和控制点。
  \end{block}
\end{frame}

\begin{frame}{例子：一个坏架构图通常坏在哪里（2）}
  \Core{学生最常见的问题是画出很多框，但无法说明数据、状态和决策谁在流动。}
\begin{block}{坏图二}
    只画数据流，不画状态驻留和资源回收。
  \end{block}
\end{frame}

\begin{frame}{例子：一个坏架构图通常坏在哪里（3）}
  \Core{学生最常见的问题是画出很多框，但无法说明数据、状态和决策谁在流动。}
\begin{block}{坏图三}
    把日志、benchmark、回归系统完全排除在架构之外，导致系统边界被画得过窄。
  \end{block}
\end{frame}

\begin{frame}{课堂练习：把架构图画成三层而不是一团（1）}
  \Core{课程作业里更推荐“三层图”：入口与控制、执行与状态、观测与验证。}
\begin{block}{练习一}
    在你的架构图中分别标出请求入口、调度器、执行器、状态管理器和指标记录点。
  \end{block}
\end{frame}

\begin{frame}{课堂练习：把架构图画成三层而不是一团（2）}
  \Core{课程作业里更推荐“三层图”：入口与控制、执行与状态、观测与验证。}
\begin{block}{练习二}
    指出哪一个模块最像“全局控制点”，哪一个模块最像“热路径执行点”。
  \end{block}
\end{frame}

\begin{frame}{课堂练习：把架构图画成三层而不是一团（3）}
  \Core{课程作业里更推荐“三层图”：入口与控制、执行与状态、观测与验证。}
\begin{block}{练习三}
    用一句话回答：为什么日志和 benchmark 不应该被排除在系统架构课之外？
  \end{block}
\end{frame}


\begin{frame}{本讲总结}
  \Core{架构视角是统一论文、代码与实验分析的关键桥梁。}
  \begin{itemize}
    \item 系统研究不能只停留在机制名词。
    \item 理解模块边界，才能设计可落地项目。
  \end{itemize}
\end{frame}

\begin{frame}{对应 Tutorial}
  \Core{本讲对应 Tutorial 6。先独立作答，再对照下一页参考答案。}
  \begin{itemize}
    \item 文件：\texttt{\detokenize{build/tutorials/Tutorial_06_推理系统架构.pdf}}
    \item 动手运行、日志记录与提交要求统一查看 \texttt{\detokenize{build/experiments/}} 和对应 experiment 源目录下的 \texttt{\detokenize{assignment_spring-2026/}}。
  \end{itemize}
\end{frame}

\end{document}